\begin{lemma}
Let $\Omega\subseteq\mathbb{C}$ be open, let $j\in\{0,1\}$, write $e_0=1$ and
$e_1=i$, and let $u,D\colon\mathbb{C}\to\mathbb{C}$ satisfy
\begin{enumerate}
\item[(H1)] $u\in L^{1}_{\mathrm{loc}}(\Omega)$ and
  $D\in L^{1}_{\mathrm{loc}}(\Omega)$, in the sense of \noderef{MemLpLocOn};
\item[(H2)] $D$ is the distributional $j$-th partial derivative of $u$ on
  $\Omega$, in the sense of \noderef{IsWeakPartialOn}.
\end{enumerate}
Let $K\subseteq\Omega$ be compact.  Then there is $\delta>0$ such that for
every real $t$ with $|t|\le\delta$,
\[
  u(z+te_j)-u(z)=\int_0^tD(z+se_j)\,ds
  \qquad\text{for almost every }z\in K .
\]

This is the absolute continuity of $u$ along almost every line parallel to
$e_j$, localised to a compact subset of $\Omega$.  It is the only route from
(H2), which sees $u$ exclusively through integrals against test functions, to
pointwise information about the increments of $u$ itself, and it is what the
paper relies on implicitly wherever Sections 4 to 6 treat the
$W^{1,2}$ derivatives of $H$, and of the inverse $\chi$, as classical
derivatives: the integration-by-parts identity alone says nothing about any
particular representative of the $L^1$-class of $u$, whereas the displayed
identity pins down every increment of $u$ along $e_j$ from $D$.  The
restriction to a compact $K$ and to small $|t|$ is not a weakening: the
translate $z+te_j$ must stay inside $\Omega$ for the right-hand side to be
meaningful, and no uniform bound on $|t|$ is available on $\Omega$ itself.
\end{lemma}
\begin{proof}
Write $d\sigma$ for plane Lebesgue measure.  Throughout, $\partial_jv(z)
=\mathrm{d}v(z)[e_j]$ denotes the classical partial derivative of a
differentiable $v$ along $e_j$.

\emph{Step 0: the scale $\delta$.}  Since $K$ is compact, $\Omega$ is open and
$K\subseteq\Omega$, there is $\delta>0$ with
\[
  K^{3\delta}:=\{z:\operatorname{dist}(z,K)\le3\delta\}\subseteq\Omega ;
\]
indeed either $\Omega=\mathbb{C}$, where any $\delta>0$ serves, or
$\operatorname{dist}(K,\mathbb{C}-\Omega)>0$ because $K$ is compact and
$\mathbb{C}-\Omega$ is closed, nonempty and disjoint from $K$, and then any
$\delta$ below a third of that distance serves.  The set $K^{3\delta}$ is
closed and bounded, hence compact, and
$K\subseteq K^{\delta}\subseteq K^{2\delta}\subseteq K^{3\delta}$.  We prove
the assertion with this $\delta$.  For $t=0$ both sides vanish, so we may
assume $0<|t|\le\delta$.

\emph{Step 1: the mollifier.}  Fix a function
$\varrho\colon\mathbb{C}\to\mathbb{R}$ of class $C^\infty$ with $\varrho\ge0$,
$\varrho(y)=0$ for $|y|\ge1$ and $\int\varrho\,d\sigma=1$, and for
$0<\varepsilon\le\delta$ put
$\varrho_\varepsilon(y)=\varepsilon^{-2}\varrho(y/\varepsilon)$, which is
$C^\infty$, vanishes for $|y|\ge\varepsilon$ and has
$\int\varrho_\varepsilon\,d\sigma=1$.  Let
\[
  U=\mathbf{1}_{K^{3\delta}}\,u,\qquad V=\mathbf{1}_{K^{3\delta}}\,D .
\]
By (H1) applied to the compact set $K^{3\delta}\subseteq\Omega$ and
\noderef{MemLpLocOn}, $U$ and $V$ belong to $L^{1}(\mathbb{C})$.  Put
$u_\varepsilon=U*\varrho_\varepsilon$ and
$v_\varepsilon=V*\varrho_\varepsilon$.

\emph{Step 2: $u_\varepsilon$ is smooth and
$\partial_ju_\varepsilon=v_\varepsilon$ near $K$.}  The convolution of an
$L^{1}$ function with a compactly supported $C^\infty$ function is $C^\infty$,
with derivatives obtained by differentiating the smooth factor; so
$u_\varepsilon$ is $C^\infty$ on $\mathbb{C}$ and
\begin{equation}\tag{1}
  \partial_ju_\varepsilon(z)
  =\bigl(U*\partial_j\varrho_\varepsilon\bigr)(z)
  =\int U(w)\,(\partial_j\varrho_\varepsilon)(z-w)\,d\sigma(w).
\end{equation}
Fix $z$ with $\operatorname{dist}(z,K)\le2\delta$ and $0<\varepsilon\le\delta$,
and define $\varphi(w)=\varrho_\varepsilon(z-w)$, read as a complex-valued
function.  Then $\varphi$ is $C^\infty$ and vanishes off the closed ball
$\{w:|w-z|\le\varepsilon\}$, which is compact and contained in
$K^{3\delta}\subseteq\Omega$ because $\operatorname{dist}(z,K)\le2\delta$ and
$\varepsilon\le\delta$.  Hence $\operatorname{tsupp}\varphi\subseteq\Omega$ and
$\varphi$ is a test function for $\Omega$ in the sense of \noderef{IsTestFn}.
The map $w\mapsto z-w$ is affine with derivative $-\mathrm{id}$, so by the
chain rule
\[
  \mathrm{d}\varphi(w)[e_j]=-(\partial_j\varrho_\varepsilon)(z-w).
\]
Therefore (H2), applied to $\varphi$, reads
\[
  -\int u(w)\,(\partial_j\varrho_\varepsilon)(z-w)\,d\sigma(w)
  =-\int D(w)\,\varrho_\varepsilon(z-w)\,d\sigma(w),
\]
that is
\begin{equation}\tag{2}
  \int u(w)\,(\partial_j\varrho_\varepsilon)(z-w)\,d\sigma(w)
  =\int D(w)\,\varrho_\varepsilon(z-w)\,d\sigma(w).
\end{equation}
Both integrands vanish off $\{w:|w-z|\le\varepsilon\}\subseteq K^{3\delta}$,
where $u=U$ and $D=V$; so by (1) the left side of (2) is
$\partial_ju_\varepsilon(z)$ and the right side is $v_\varepsilon(z)$.  Hence
\begin{equation}\tag{3}
  \partial_ju_\varepsilon(z)=v_\varepsilon(z)
  \qquad\bigl(\operatorname{dist}(z,K)\le2\delta,\ 0<\varepsilon\le\delta\bigr).
\end{equation}

\emph{Step 3: the finite-difference identity for $u_\varepsilon$.}  Let
$z\in K$ and let $s$ lie between $0$ and $t$.  Then
$\operatorname{dist}(z+se_j,K)\le|s|\le|t|\le\delta\le2\delta$, so (3) applies
at $z+se_j$.  The function $s\mapsto u_\varepsilon(z+se_j)$ is $C^\infty$ on
$\mathbb{R}$ with derivative $\partial_ju_\varepsilon(z+se_j)$, by the chain
rule, so the fundamental theorem of calculus and (3) give
\begin{equation}\tag{4}
  u_\varepsilon(z+te_j)-u_\varepsilon(z)
  =\int_0^t\partial_ju_\varepsilon(z+se_j)\,ds
  =\int_0^tv_\varepsilon(z+se_j)\,ds .
\end{equation}

\emph{Step 4: letting $\varepsilon\to0$ in (4).}  For $g\in L^{1}(\mathbb{C})$
the mollifications converge in $L^1$, that is
$\|g*\varrho_\varepsilon-g\|_{L^{1}(\mathbb{C})}\to0$ as $\varepsilon\to0$, by
\noderef{MollificationL1Convergence}; applied to $U$ and $V$,
\begin{equation}\tag{5}
  \|u_\varepsilon-U\|_{L^{1}(\mathbb{C})}\to0,\qquad
  \|v_\varepsilon-V\|_{L^{1}(\mathbb{C})}\to0 .
\end{equation}
Regard both sides of (4) as functions of $z\in K$, with $t$ fixed.  For the
left side, translation invariance of $d\sigma$ gives
\[
  \int_K\bigl|u_\varepsilon(z+te_j)-U(z+te_j)\bigr|\,d\sigma(z)
  \le\|u_\varepsilon-U\|_{L^{1}(\mathbb{C})},
\]
and the same bound without the translation, so
$z\mapsto u_\varepsilon(z+te_j)-u_\varepsilon(z)$ converges in $L^{1}(K)$ to
$z\mapsto U(z+te_j)-U(z)$, which equals $u(z+te_j)-u(z)$ for $z\in K$ because
$z$ and $z+te_j$ both lie in $K^{3\delta}$.  For the right side, write $I_t$
for the interval between $0$ and $t$; the function
$(z,s)\mapsto\bigl(v_\varepsilon-V\bigr)(z+se_j)$ is measurable on
$K\times I_t$, being the composition of a measurable function with the
continuous map $(z,s)\mapsto z+se_j$, so by Tonelli's theorem and translation
invariance
\[
  \int_K\Bigl|\int_0^t\bigl(v_\varepsilon-V\bigr)(z+se_j)\,ds\Bigr|
    d\sigma(z)
  \le\int_{I_t}\!\Bigl(\int_K\bigl|v_\varepsilon-V\bigr|(z+se_j)
    \,d\sigma(z)\Bigr)ds
  \le|t|\,\|v_\varepsilon-V\|_{L^{1}(\mathbb{C})},
\]
which tends to $0$ by (5); and $V(z+se_j)=D(z+se_j)$ for $z\in K$ and
$|s|\le\delta$, because then $z+se_j\in K^{\delta}\subseteq K^{3\delta}$.
Thus the left sides of (4) converge in $L^1(K)$ to
$z\mapsto u(z+te_j)-u(z)$ and the right sides converge in $L^1(K)$ to
$z\mapsto\int_0^tD(z+se_j)\,ds$.  Two $L^{1}(K)$ limits of one and the same
family agree almost everywhere, so
\[
  u(z+te_j)-u(z)=\int_0^tD(z+se_j)\,ds
  \qquad\text{for almost every }z\in K,
\]
which is the assertion for this $t$.  As $t$ with $0<|t|\le\delta$ was
arbitrary, and the case $t=0$ is trivial, the proof is complete.
\end{proof}
